% ==============================================================================
% Page 2: Abstract
% ==============================================================================
\newpage
\pagenumbering{arabic}
\renewcommand{\thepage}{(1)}

\begin{center}
    {\fontsizeSubtitle\bfseries Abstract \par}
\end{center}
\vspace{1em}

\noindent
This report describes my cooperative education program at \companyName{}, a game development company, from 25 August to 30 November 2026. I developed a personal 2D top-down action game in Unity and followed the phases of a professional workflow: pitching, prototyping, pre-production documents (a game design document, a technical design document, and a UML class diagram), and production with pull request review.
The report asks how this professional workflow differs from my earlier solo workflow, and what each one costs and gives.
I compare the two using time data from the company project and from Mashima, a game I developed alone during my previous cooperative education program in Japan.
The professional workflow was slower: pre-production took about half of the internship, and the same player movement that took about 7 hours in the prototype took about 3.4 times as long in production.
In return, the design was shared in documents, and the code was reviewed before it was merged.
The scope of the game was reduced to a single Time Survival stage so that it could be delivered as a complete game.
I conclude that the professional workflow is easier to scale and maintain, while the solo workflow is faster for exploring an idea.
The internship also showed me that planning, communication, and asking questions early are the skills I most need to improve.